% ══════════════════════════════════════════════════════
%  contenu/page-de-garde-exponentielle.tex
%  Page de garde - Fonction Exponentielle
% ══════════════════════════════════════════════════════
\thispagestyle{empty}
\usetikzlibrary{intersections,decorations.text}

\begin{tikzpicture}[overlay,remember picture,font=\sffamily\bfseries]

  %% ─── Arc principal ────────────────────────────────
  \draw[very thick,c4,name path=bigarc]
    ([xshift=-2mm]current page.north) arc(150:285:11)
    coordinate[pos=0.22](x0);

  \begin{scope}
    \clip ([xshift=-2mm]current page.north) arc(150:285:11)
      -- (current page.north east);
    \fill[c4!40,opacity=0.3] ([xshift=4.5cm]x0) circle(4.5);
    \fill[c4!40,opacity=0.3] ([xshift=3.3cm]x0) circle(3.3);
    \fill[c4!40,opacity=0.3] ([xshift=2.2cm]x0) circle(2.2);
    \draw[very thick,c4!65] (x0) arc(-90:30:6.5);
    \draw[very thick,c4]    (x0) arc(90:-30:8.75);
    \draw[very thick,c4!65,name path=arc1] (x0) arc(90:-90:4.65);
    \draw[very thick,c4!50] (x0) arc(90:-90:2.85);
    \path[name intersections={of=bigarc and arc1,by=x1}];
    \draw[very thick,c4,name path=arc2] (x1) arc(135:-20:4.75);
    \draw[very thick,c4!55] (x1) arc(135:-20:8.75);
    \path[name intersections={of=bigarc and arc2,by={aux,x2}}];
    \draw[very thick,c4!50] (x2) arc(180:50:2.2);
  \end{scope}

  %% ─── Année sur l'arc ──────────────────────────────
  \path[decoration={text along path,text color=c4,
      raise=-2.8ex,text along path,
      text={|\sffamily\bfseries|2025 --- 2026},
      text align=center},decorate]
    ([xshift=-2mm]current page.north) arc(150:245:11);

  %% ─── Bloc décoratif central droit ────────────────
  \begin{scope}
    \path[clip,postaction={fill=c3!80}]
      ([xshift=2.2cm,yshift=-7.8cm]current page.center)
      rectangle ++(4.5,7.8);
    \fill[c4!90]
      ([xshift=0.5cm,yshift=-7.8cm]current page.center)
      arc(180:60:2) |- ++(-3,6.2) --cycle;
    \foreach \dx in {-1.5,0.5,2.5,4.5}{
      \draw[very thick,c4!40]
        ([xshift=\dx cm,yshift=-7.8cm]current page.center)
        arc(180:0:2);}
    %% Formule spécifique à l'exponentielle
    \node[text=white,scale=1.3] at
      ([xshift=4.8cm,yshift=-4.5cm]current page.center)
      {$e^x=\displaystyle\sum_{n=0}^{\infty}\frac{x^n}{n!}$};
  \end{scope}

  %% ─── Fond principal gauche ────────────────────────
  \fill[c1]
    ([xshift=2.2cm,yshift=-7.8cm]current page.center)
    rectangle ++(-14.5,7.8);

  %% ─── Textes principaux ───────────────────────────
  \node[text=white,anchor=west,
        font=\bfseries\fontsize{38}{44}\sffamily]
    at ([xshift=-9.5cm,yshift=-2.2cm]current page.center)
    {Math\'ematiques};

  \node[text=white!80,anchor=west,
        font=\bfseries\fontsize{18}{22}\sffamily]
    at ([xshift=-9.5cm,yshift=-3.9cm]current page.center)
    {2\textsuperscript{\`eme} Bac Sciences Math\'ematiques};

  \node[text=HeaderOrange,anchor=west,
        font=\bfseries\fontsize{22}{28}\sffamily]
    at ([xshift=-9.5cm,yshift=-5.5cm]current page.center)
    {Chapitre 5 : Fonction Exponentielle};

  \node[text=white!60,anchor=west,
        font=\fontsize{13}{16}\sffamily]
    at ([xshift=-9.5cm,yshift=-6.8cm]current page.center)
    {Cours \textbullet\ Exercices r\'esolus \textbullet\
     Probl\`emes};

  %% ─── Graphique de la fonction exponentielle ──────
  \begin{scope}[xshift=4cm,yshift=-3cm]
    \draw[->] (-2,0) -- (2.5,0) node[right] {$x$};
    \draw[->] (0,-0.5) -- (0,3.5) node[above] {$y$};
    \draw[domain=-2:1.5, smooth, thick, MainRed]
      plot (\x, {exp(\x)});
    \node[right, MainRed] at (1.5,4.5) {$y=e^x$};
  \end{scope}

  %% ─── Cadre général ────────────────────────────────
  \draw[c4!50,line width=4pt]
    ([xshift=3mm,yshift=-3mm]current page.south west)
    rectangle
    ([xshift=-3mm,yshift=3mm]current page.north east);

  %% ─── Auteur ──────────────────────────────────────
  \node[text=gray!55,font=\itshape\large\sffamily,
        anchor=south]
    at ([yshift=2cm]current page.south)
    {Pr\'epar\'e par : A.\ EL ROMANI
     \quad\textemdash\quad 2025--2026};

\end{tikzpicture}
\vfill\newpage